\subsection{Problem Formulation}
Consider a federated graph learning setting with $K$ clients. Each client $k \in [K]$ holds a local graph $\mathcal{G}_k = (\mathcal{V}_k, \mathcal{E}_k, \mathbf{X}_k)$, where $\mathcal{V}_k$ is the set of nodes, $\mathcal{E}_k$ is the set of edges, and $\mathbf{X}_k \in \mathbb{R}^{|\mathcal{V}_k| \times d}$ is the feature matrix. Each node $v_i \in \mathcal{V}_k$ has a label $y_i \in \{0, 1\}$ (0 for legitimate, 1 for fraud) and a sensitive attribute $s_i \in \{0, 1\}$.

The goal is to collaboratively train a GNN model $f_\theta$ to minimize the global loss function $\mathcal{L}(\theta)$ while satisfying a global fairness constraint and preserving differential privacy. We define the fairness metric as the Demographic Parity Difference (DPD):
\begin{equation}
\Delta_{DP} = |P(\hat{y}=1 | s=0) - P(\hat{y}=1 | s=1)| \leq \tau
\end{equation}
where $\tau$ is a pre-defined fairness budget.

\subsection{Fairness-Sensitive Edge Reweighting (FSER)}
Graph homophily often leads to biased information propagation. FSER aims to mitigate this by reweighting edges during the aggregation phase. We adopt a Graph Attention Network (GAT) backbone. The standard attention coefficient $e_{ij}$ between node $i$ and neighbor $j$ is computed as:
\begin{equation}
e_{ij} = \text{LeakyReLU}(\mathbf{a}^T [\mathbf{W}\mathbf{x}_i || \mathbf{W}\mathbf{x}_j])
\end{equation}

To incorporate fairness, we define a \textbf{Fairness Risk Score} $\phi_{ij}$ that penalizes edges between nodes with different sensitive attributes if their feature embeddings are highly similar (indicating potential confounding):
\begin{equation}
\phi_{ij} = \mathbb{I}(s_i \neq s_j) \cdot \text{ReLU}(\cos(\mathbf{x}_i, \mathbf{x}_j))
\end{equation}
The reweighted attention coefficients are then:
\begin{equation}
\alpha_{ij} = \frac{\exp(e_{ij} - \beta \phi_{ij})}{\sum_{k \in \mathcal{N}_i} \exp(e_{ik} - \beta \phi_{ik})}
\end{equation}
where $\beta$ is a hyperparameter controlling the strength of the fairness penalty. This mechanism effectively down-weights edges that are likely to propagate biased information.

\subsection{Fairness-Task Gradient Decomposition (FTGD)}
Standard Differential Privacy (DP) adds noise to the entire gradient, often degrading utility. We propose FTGD to decouple task-specific information from fairness-sensitive information. Let $\nabla \mathcal{L}_{task}$ be the gradient of the utility loss (e.g., BCE) and $\nabla \mathcal{L}_{fair}$ be the gradient of the local fairness regularization term.

We decompose the total gradient $g = \nabla \mathcal{L}_{task} + \lambda \nabla \mathcal{L}_{fair}$ into two orthogonal components:
\begin{align}
g_{fair} &= \nabla \mathcal{L}_{fair} \\
g_{task}^{\perp} &= g - \text{proj}_{g_{fair}}(g) = g - \frac{\langle g, g_{fair} \rangle}{\|g_{fair}\|^2} g_{fair}
\end{align}
Since sensitive information is primarily concentrated in $g_{fair}$, we apply Gaussian noise only to this component:
\begin{equation}
\tilde{g}_{fair} = Clip(g_{fair}, C) + \mathcal{N}(0, \sigma^2 \mathbf{I})
\end{equation}
where $C$ is the clipping threshold and $\sigma$ is calibrated to satisfy $(\epsilon, \delta)$-DP. The final update sent to the server is $g_{update} = g_{task}^{\perp} + \tilde{g}_{fair}$. This preserves the utility-critical direction $g_{task}^{\perp}$ without noise.

\subsection{Bi-Objective Frank-Wolfe Aggregation (BFWA)}
The server aggregates client updates to maximize utility subject to the global fairness constraint. This is formulated as a constrained optimization problem:
\begin{align}
\min_{\mathbf{w}} \quad & \sum_{k=1}^K w_k \mathcal{L}_{k}(\theta) \\
\text{s.t.} \quad & \sum_{k=1}^K w_k \Delta_{DP, k}(\theta) \leq \tau, \quad \sum w_k = 1, w_k \geq 0
\end{align}
To solve this efficiently, we employ the Frank-Wolfe algorithm. We define the Lagrangian $\mathcal{L}(\mathbf{w}, \mu) = f(\mathbf{w}) + \mu (g(\mathbf{w}) - \tau)$. At each round $t$, we compute the gradient $\nabla_{\mathbf{w}} \mathcal{L}$ and find the linear minimization oracle (LMO) solution $\mathbf{s}_t$:
\begin{equation}
\mathbf{s}_t = \arg\min_{\mathbf{s} \in \Delta_K} \langle \nabla_{\mathbf{w}} \mathcal{L}(\mathbf{w}_t, \mu_t), \mathbf{s} \rangle
\end{equation}
The weights are updated via $\mathbf{w}_{t+1} = (1-\gamma_t)\mathbf{w}_t + \gamma_t \mathbf{s}_t$, and the dual variable $\mu$ is updated via gradient ascent:
\begin{equation}
\mu_{t+1} = \max(0, \mu_t + \eta_\mu (\sum w_{k,t} \Delta_{DP,k} - \tau))
\end{equation}
This allows the server to dynamically upweight clients that contribute to fairness while maintaining high utility.

\begin{algorithm}
\caption{Bi-Objective Frank-Wolfe Aggregation (BFWA)}
\begin{algorithmic}[1]
\STATE \textbf{Input:} Client gradients $\{g_k\}$, metrics $\{perf_k, dpd_k\}$, budget $\tau$
\STATE Initialize weights $\mathbf{w} = [1/K, \dots, 1/K]$
\STATE Initialize dual variable $\mu = 0$
\FOR{$t = 0 \to T_{FW}$}
    \STATE Compute gradient $\nabla \mathcal{L} = -Perf + \mu \cdot DPD$
    \STATE Find LMO: $k^* = \arg\min \nabla \mathcal{L}$
    \STATE $\mathbf{s} = \mathbf{0}$; $s_{k^*} = 1$
    \STATE $\mathbf{w} \leftarrow (1-\gamma_t)\mathbf{w} + \gamma_t \mathbf{s}$
    \STATE Update $\mu \leftarrow \max(0, \mu + \eta (\mathbf{w}^T \mathbf{DPD} - \tau))$
\ENDFOR
\RETURN $\mathbf{w}$
\end{algorithmic}
\label{alg:bfwa}
\end{algorithm}
